\section{Related Work}
\label{sec:rw}

IoT middleware standards such as OneM2M provide interoperability across vendors and device types, but leave dynamic Quality-of-Service (QoS) management and traffic-offload decisions largely unaddressed, forcing operators to rely on manual, reactive intervention when systems become overloaded.

\subsection{Baseline: ML/DL for OneM2M QoS Optimization}
The baseline for this project is Et-Tousy, Zyane \& Sharif \cite{EtTousy26}, ``Adaptive QoS Management in OneM2M Standard: Machine Learning and Deep Learning for IoT Network Optimization'' (Journal of Network and Systems Management, May 2026). The paper collects real-time QoS telemetry (RTT, CPU, RAM, success rate) from Azure IoT devices via a MAPE-K monitoring loop, labels each window as Preferable / Acceptable / Critical using fixed thresholds, and trains classical ML models (Random Forest, SVC, KNN, XGBoost, LightGBM), ensembles (Voting/Stacking of RF+LightGBM), and deep learning models (CNN, LSTM, VAE+RF) to recommend a traffic-offload decision: keep processing local, partially offload to the cloud, or fully offload. Their best model (a Random Forest / LightGBM ensemble) reaches 98\% accuracy and is deployed as a live API that drives the Analyzer stage of a MAPE-K loop, reducing RTT by 30-50\% and increasing success rates by 10-15\% across uniform, burst, and real-time traffic scenarios.

\subsection{The Gap This Project Targets}
The paper is explicit about what it does not address. In their own Discussion section: ``the system currently relies on a centralized learning mechanism, which might not scale well or guarantee data privacy in distributed IoT environments. To overcome this, we plan to explore federated learning.'' This is stated as future work, not built or evaluated in the paper itself.

\subsection{Federated Learning for Edge/IoT Systems}
Federated learning is a well-established response to exactly this concern -- training models across distributed clients without centralizing raw data -- but the baseline paper's own centralized Random Forest pipeline does not use it. This project reproduces that centralized baseline, then builds and evaluates a federated alternative: one Random Forest per simulated IoT site, trained only on that site's local telemetry, combined at inference time by averaging predicted class probabilities across sites. No raw telemetry is pooled across sites at any point, directly addressing the privacy and centralization concern the baseline paper names but does not resolve.
